% !PRINTED 232
\weberpage{232}
%<<CONT
setzen und mit $h$, $k$ zwei Constanten bezeichnen, so ist
\[
\Phi(y) = h\,\Phi(x), \quad \Psi(y) = k\,\Psi(x),
\]
vorausgesetzt, dass $\begin{pmatrix} \alpha, & \beta \\ \gamma, & \delta \end{pmatrix}$ eine der Ikosa\"edersubstitutionen ist. Da nun aber $y$ sowohl als $x$ eine unabhängige Variable ist, so besteht nach (7) die Identität
\[
\Phi(x) + \Psi(x) = h\,\Phi(x) + k\,\Psi(x) = 1728,
\]
und da $\Phi$ und $\Psi$ nicht in constantem Verhältnisse stehen, so muss $h = k = 1$ sein, d. h. die Quotienten $\Phi$ und $\Psi$ bleiben bei den Ikosa\"edersubstitutionen absolut ungeändert.

Daraus ergiebt sich weiter, \gesperrt{dass $f$, $H$, $T$ bei den homogenen Ikosa\"edersubstitutionen mit der Determinante $1$ absolut ungeändert bleiben.}

Denn wenn $f$ durch eine solche Substitution in $cf$ übergeht, so gehen $H$ und $T$ in $c^2H$ und $c^3T$ über [nach (2) und (4)], und die Quotienten $\Phi$ und $\Psi$ in $c\Phi$ und $c\Psi$. Da andererseits $\Phi$ und $\Psi$ ungeändert bleiben, so ist $c = 1$.

Setzen wir $\Psi(x) = z$, so wird $\Phi(x) = 1728 - z$, und wir erhalten aus (8) die beiden Gleichungen:
\weq{9}{H^3 - z f^5 = 0,}
\weq{10}{T^2 - (1728 - z)\,f^5 = 0,}
von denen wegen der Identität (6) die eine aus der anderen folgt, so dass es eigentlich nur zwei verschiedene Formen für eine und dieselbe Gleichung sind. Betrachten wir nun darin $z$ als gegeben, so haben wir eine Gleichung 60\textsuperscript{sten} Grades für $x$, die die \gesperrt{Ikosa\"edergleichung} heisst.

Auf die Eigenschaften dieser Gleichung und ihre Beziehung zu der allgemeinen Gleichung 5\textsuperscript{ten} Grades kommen wir in einem späteren Abschnitte zurück.

\begin{center}
\S.\,61.\\[2pt]
\gesperrt{Die Invarianten des Ikosa\"eders.}
\end{center}

Durch die Grundformen der Poly\"edergruppen, die wir bisher kennen gelernt haben, ist, wie wir jetzt beweisen können, das Gebiet der Invarianten der betreffenden Gruppen erschöpft. Wir
%CONT>>
